% =========================================================
% Proof Grep Stubs — Rational Numbers (Chapter 10)
% Source: Learning_Real_Analysis (76).pdf
% Purpose: Combined proof-stub file for theorem-like statements
% Mode: standard stubs only (no completed proofs)
% =========================================================

\section*{Proof Grep — Rational Numbers}

This file collects proof stubs for the theorem-like statements in Chapter~10
of the current Rational Numbers PDF.  The statements have been arranged in a
\emph{dependency-aware order} so that earlier stubs do not rely on later ones.

\begin{remark*}[Dependency order used]
The proof stubs below are ordered so that structural results about the quotient
construction and ordered-field consequences appear before density, mediants,
Farey phenomena, rational approximation, Cauchy-sequence results, and finally
limiting-process results involving Stern--Brocot refinement.
\end{remark*}

\begin{remark*}[Duplicate formulations consolidated]
The current PDF contains three duplicate theorem/corollary pairs:
\begin{itemize}
  \item Theorem~10.25 and Theorem~10.40 both state $1 > 0$.
  \item Corollary~10.26 and Corollary~10.41 both state $n > 0$ for all $n \in \mathbb{N}$.
  \item Theorem~10.33 and Theorem~10.49 are both mediant inequalities.
\end{itemize}
To avoid redundant stubs, this file keeps the canonical occurrences
Theorem~10.25, Corollary~10.26, and Theorem~10.49.
\end{remark*}

\subsection*{Theorem 10.6 — The relation $\sim$ is an equivalence relation on $W$}
\label{prf:grep-rational-equivalence-relation}

\begin{proof}
\Claim The relation $\sim$ on $W$ defined by
\[
(a,b) \sim (c,d) \iff ad = bc
\]
is an equivalence relation on $W$.

\textbf{Given.} The set $W = \mathbb{Z} \times (\mathbb{Z}\setminus\{0\})$ and the relation above.

\textbf{Goal.} Show that $\sim$ is reflexive, symmetric, and transitive.

\textbf{Strategy.} Verify the three defining properties of an equivalence relation directly from the defining equation $ad=bc$.

% -------------------------------------------------------
% YOUR PROOF HERE
% -------------------------------------------------------

\end{proof}

\begin{remark*}[Proof shape]
This is a direct verification proof.  The reflexive and symmetric parts are
immediate algebraic identities; the transitive part is the only nontrivial step
and should be handled by chaining two cross-multiplication equalities.
\end{remark*}

\begin{remark*}[Dependencies]
Uses only the definition of $W$, the definition of $\sim$, and basic arithmetic
of the integers.
\end{remark*}


\subsection*{Theorem 10.21 — For every $a \in F$, $a^2 \ge 0$}
\label{prf:grep-squares-nonnegative}

\begin{proof}
\Claim For every $a \in F$,
\[
a^2 \ge 0.
\]

\textbf{Given.} An ordered field $F$ presented via a positive cone $P$.

\textbf{Goal.} Show that every square is positive or zero.

\textbf{Strategy.} Apply trichotomy to $a$ and split into the cases $a \in P$, $a=0$, and $-a \in P$.

% -------------------------------------------------------
% YOUR PROOF HERE
% -------------------------------------------------------

\end{proof}

\begin{remark*}[Proof shape]
This is a three-case proof obtained by unpacking the positive-cone axioms.
The key observation is that if $-a$ is positive, then $(-a)(-a)=a^2$.
\end{remark*}

\begin{remark*}[Dependencies]
Uses Definition~10.19 (positive cone), especially trichotomy and closure under
multiplication.
\end{remark*}


\subsection*{Theorem 10.25 — In an ordered field, $1 > 0$}
\label{prf:grep-one-positive}

\begin{proof}
\Claim In an ordered field,
\[
1 > 0.
\]

\textbf{Given.} An ordered field (or equivalently a field with a positive cone).

\textbf{Goal.} Show that the multiplicative identity is positive.

\textbf{Strategy.} Use trichotomy on $1$.  Eliminate the possibilities $1=0$ and $-1>0$.

% -------------------------------------------------------
% YOUR PROOF HERE
% -------------------------------------------------------

\end{proof}

\begin{remark*}[Proof shape]
This is a contradiction-by-cases proof.  The case $1=0$ contradicts the field
axioms, and the case $-1>0$ collapses under multiplication.
\end{remark*}

\begin{remark*}[Dependencies]
Uses Definition~10.19 (positive cone), the field axiom $1 \neq 0$, and closure
of positive elements under multiplication.
\end{remark*}


\subsection*{Theorem 10.22 — If $a>0$, then $a^{-1}>0$}
\label{prf:grep-positive-inverse}

\begin{proof}
\Claim If $a>0$, then
\[
a^{-1} > 0.
\]

\textbf{Given.} An ordered field and an element $a$ with $a>0$.

\textbf{Goal.} Show that the multiplicative inverse of $a$ is positive.

\textbf{Strategy.} Suppose the inverse were not positive and derive a contradiction by multiplying by a positive element.

% -------------------------------------------------------
% YOUR PROOF HERE
% -------------------------------------------------------

\end{proof}

\begin{remark*}[Proof shape]
This is a contradiction proof.  The expected contradiction is that a positive
number multiplied by an allegedly negative inverse would force $-1$ to be
positive, contradicting Theorem~10.25.
\end{remark*}

\begin{remark*}[Dependencies]
Uses Theorem~10.25, Definition~10.19, and the field axiom that every nonzero
element has a multiplicative inverse.
\end{remark*}


\subsection*{Theorem 10.23 — If $a<b$ and $c>0$, then $ac<bc$}
\label{prf:grep-multiply-by-positive}

\begin{proof}
\Claim If $a<b$ and $c>0$, then
\[
ac < bc.
\]

\textbf{Given.} Elements $a,b,c$ in an ordered field with $a<b$ and $c>0$.

\textbf{Goal.} Show that multiplication by a positive element preserves strict inequality.

\textbf{Strategy.} Translate $a<b$ into positivity of $b-a$, then multiply by $c$ and translate back.

% -------------------------------------------------------
% YOUR PROOF HERE
% -------------------------------------------------------

\end{proof}

\begin{remark*}[Proof shape]
This is a direct proof entirely inside the positive-cone language:
$b-a \in P$, $c \in P$, hence $c(b-a) \in P$.
\end{remark*}

\begin{remark*}[Dependencies]
Uses Definition~10.19 and closure of the positive cone under multiplication.
\end{remark*}


\subsection*{Theorem 10.24 — If $a<b$ and $c<0$, then $ac>bc$}
\label{prf:grep-multiply-by-negative}

\begin{proof}
\Claim If $a<b$ and $c<0$, then
\[
ac > bc.
\]

\textbf{Given.} Elements $a,b,c$ in an ordered field with $a<b$ and $c<0$.

\textbf{Goal.} Show that multiplication by a negative element reverses strict inequality.

\textbf{Strategy.} Replace $c$ by $-c>0$, apply Theorem~10.23, and then rewrite the resulting inequality.

% -------------------------------------------------------
% YOUR PROOF HERE
% -------------------------------------------------------

\end{proof}

\begin{remark*}[Proof shape]
This is a reduction proof: first convert the negative multiplier into a
positive one, then appeal to the previous theorem.
\end{remark*}

\begin{remark*}[Dependencies]
Uses Theorem~10.23 and Definition~10.19.
\end{remark*}


\subsection*{Corollary 10.26 — For every natural number $n$, $n>0$}
\label{prf:grep-natural-numbers-positive}

\begin{proof}
\Claim For every natural number $n$,
\[
n > 0.
\]

\textbf{Given.} An ordered field and a natural number $n$ regarded as a finite sum of copies of $1$.

\textbf{Goal.} Show that every natural number is positive.

\textbf{Strategy.} Combine Theorem~10.25 with closure of the positive cone under addition.

% -------------------------------------------------------
% YOUR PROOF HERE
% -------------------------------------------------------

\end{proof}

\begin{remark*}[Proof shape]
This is a short consequence proof.  Once $1>0$ is known, every finite sum of
copies of $1$ is positive.
\end{remark*}

\begin{remark*}[Dependencies]
Uses Theorem~10.25 and Definition~10.19.
\end{remark*}


\subsection*{Theorem 10.42 — Archimedean property of $\mathbb{Q}$}
\label{prf:grep-archimedean-q}

\begin{proof}
\Claim For every $x \in \mathbb{Q}$ there exists $n \in \mathbb{N}$ such that
\[
n > x.
\]
Equivalently, for every $x>0$ there exists $n \in \mathbb{N}$ such that
\[
\frac{1}{n} < x.
\]

\textbf{Given.} A rational number $x$.

\textbf{Goal.} Produce a natural number that exceeds $x$.

\textbf{Strategy.} Write $x$ in fraction form and choose a natural number large enough relative to its numerator and denominator.

% -------------------------------------------------------
% YOUR PROOF HERE
% -------------------------------------------------------

\end{proof}

\begin{remark*}[Proof shape]
This is a construction proof specific to $\mathbb{Q}$.  The natural approach is
to write $x=a/b$ and choose an explicit natural number in terms of $a$ and $b$.
\end{remark*}

\begin{remark*}[Dependencies]
Uses the construction of $\mathbb{Q}$, Corollary~10.26, and basic arithmetic of
integers and rational numbers.
\end{remark*}


\subsection*{Theorem 10.44 — Density of $\mathbb{Q}$ in $\mathbb{R}$}
\label{prf:grep-density-q-in-r}

\begin{proof}
\Claim If $x,y \in \mathbb{R}$ and $x<y$, then there exists $q \in \mathbb{Q}$ such that
\[
x < q < y.
\]

\textbf{Given.} Real numbers $x<y$.

\textbf{Goal.} Find a rational number strictly between them.

\textbf{Strategy.} Use the Archimedean property to choose $n$ with $1/n < y-x$, then locate an integer $m$ with $x < m/n \le x+1/n$.

% -------------------------------------------------------
% YOUR PROOF HERE
% -------------------------------------------------------

\end{proof}

\begin{remark*}[Proof shape]
This is the standard Archimedean-density argument.  The key move is to pass
from interval length to a denominator and then select the least integer above
$nx$.
\end{remark*}

\begin{remark*}[Dependencies]
Uses Theorem~10.42 and the order structure of the real numbers.
\end{remark*}


\subsection*{Corollary 10.45 — Between any two distinct real numbers there are infinitely many rationals}
\label{prf:grep-infinitely-many-rationals}

\begin{proof}
\Claim Between any two distinct real numbers there are infinitely many rational numbers.

\textbf{Given.} Real numbers $x<y$.

\textbf{Goal.} Show that the interval $(x,y)$ contains infinitely many rationals.

\textbf{Strategy.} Apply Theorem~10.44 repeatedly to smaller subintervals.

% -------------------------------------------------------
% YOUR PROOF HERE
% -------------------------------------------------------

\end{proof}

\begin{remark*}[Proof shape]
This is an iteration proof.  Once one rational is found, the density theorem can
be re-applied to the left and right subintervals or to a recursively chosen chain
of subintervals.
\end{remark*}

\begin{remark*}[Dependencies]
Uses Theorem~10.44.
\end{remark*}


\subsection*{Theorem 10.49 — Mediant inequality}
\label{prf:grep-mediant-inequality}

\begin{proof}
\Claim If
\[
\frac{a}{b} < \frac{c}{d}
\quad\text{with } b,d>0,
\]
then
\[
\frac{a}{b} < \frac{a+c}{b+d} < \frac{c}{d}.
\]

\textbf{Given.} Two rational numbers with positive denominators and strict order.

\textbf{Goal.} Show that the mediant lies strictly between them.

\textbf{Strategy.} Start from the cross-multiplied inequality $ad<bc$ and manipulate it separately to obtain the left-hand and right-hand bounds.

% -------------------------------------------------------
% YOUR PROOF HERE
% -------------------------------------------------------

\end{proof}

\begin{remark*}[Proof shape]
This is a two-inequality algebra proof.  Each side comes from adding a suitable
common term to $ad<bc$.
\end{remark*}

\begin{remark*}[Dependencies]
Uses only the order properties of rational numbers and basic algebraic
manipulation.
\end{remark*}


\subsection*{Corollary 10.50 — Density of $\mathbb{Q}$ within itself}
\label{prf:grep-density-q-within-q}

\begin{proof}
\Claim If $r,s \in \mathbb{Q}$ with $r<s$, then there exists $t \in \mathbb{Q}$ such that
\[
r < t < s.
\]

\textbf{Given.} Rational numbers $r<s$.

\textbf{Goal.} Produce a rational number strictly between them.

\textbf{Strategy.} Write $r$ and $s$ with positive denominators and apply the mediant inequality.

% -------------------------------------------------------
% YOUR PROOF HERE
% -------------------------------------------------------

\end{proof}

\begin{remark*}[Proof shape]
This is the immediate corollary of Theorem~10.49.
\end{remark*}

\begin{remark*}[Dependencies]
Uses Theorem~10.49 and the fact that rational numbers can be represented with
positive denominators.
\end{remark*}


\subsection*{Theorem 10.61 — Farey Neighbor Theorem}
\label{prf:grep-farey-neighbor-theorem}

\begin{proof}
\Claim Let
\[
\frac{a}{b} < \frac{c}{d}
\]
be two consecutive fractions in the Farey sequence $F_n$. Then
\[
bc-ad = 1.
\]

\textbf{Given.} Consecutive fractions in a Farey sequence.

\textbf{Goal.} Show that their determinant difference is exactly $1$.

\textbf{Strategy.} Assume the determinant difference is at least $2$, then construct a reduced fraction with denominator at most $n$ lying strictly between them, contradicting consecutiveness.

% -------------------------------------------------------
% YOUR PROOF HERE
% -------------------------------------------------------

\end{proof}

\begin{remark*}[Proof shape]
This is a contradiction proof.  The heart of the argument is that if the gap is
``too large'' arithmetically, then an intermediate Farey fraction must exist.
\end{remark*}

\begin{remark*}[Dependencies]
Uses the definition of the Farey sequence and elementary properties of reduced
fractions and inequalities.
\end{remark*}


\subsection*{Theorem 10.65 — Gap Between Farey Neighbors}
\label{prf:grep-farey-gap}

\begin{proof}
\Claim Let
\[
\frac{a}{b} < \frac{c}{d}
\]
be consecutive fractions in a Farey sequence. Then
\[
\frac{c}{d} - \frac{a}{b} = \frac{1}{bd}.
\]

\textbf{Given.} Consecutive Farey neighbors.

\textbf{Goal.} Compute their exact spacing.

\textbf{Strategy.} Subtract the two fractions and then use the determinant identity from Theorem~10.61.

% -------------------------------------------------------
% YOUR PROOF HERE
% -------------------------------------------------------

\end{proof}

\begin{remark*}[Proof shape]
This is a one-line computation once the Farey determinant condition is known.
\end{remark*}

\begin{remark*}[Dependencies]
Uses Theorem~10.61 and direct fraction subtraction.
\end{remark*}


\subsection*{Theorem 10.52 — Farey neighbor determinant condition}
\label{prf:grep-farey-neighbor-determinant-condition}

\begin{proof}
\Claim Suppose
\[
\frac{a}{b} < \frac{c}{d}
\quad\text{with } \gcd(a,b)=\gcd(c,d)=1,\; b,d>0.
\]
If
\[
bc-ad = 1,
\]
then the two fractions are neighbors in a Farey sequence, and their mediant
\[
\frac{a+c}{b+d}
\]
is already in reduced form.

\textbf{Given.} Two reduced fractions with positive denominators satisfying $bc-ad=1$.

\textbf{Goal.} Prove both the Farey-neighbor conclusion and the reduced-form conclusion for the mediant.

\textbf{Strategy.} For the neighbor claim, argue that any hypothetical reduced fraction strictly between them would force a contradiction with the determinant condition.  For reducedness of the mediant, show that a common divisor would have to divide $bc-ad=1$.

% -------------------------------------------------------
% YOUR PROOF HERE
% -------------------------------------------------------

\end{proof}

\begin{remark*}[Proof shape]
This is a two-part proof: a best-approximation/neighbor argument and an
arithmetic gcd argument.
\end{remark*}

\begin{remark*}[Dependencies]
Uses the definition of Farey sequence, basic gcd properties, and the arithmetic
identity $bc-ad=1$.
\end{remark*}


\subsection*{Theorem 10.62 — Mediant Property of Farey Neighbors}
\label{prf:grep-mediant-farey-neighbors}

\begin{proof}
\Claim If
\[
\frac{a}{b},\frac{c}{d}
\]
are consecutive fractions in a Farey sequence, then the mediant
\[
\frac{a+c}{b+d}
\]
lies strictly between them.

\textbf{Given.} Consecutive Farey neighbors.

\textbf{Goal.} Show that their mediant sits strictly between them.

\textbf{Strategy.} Apply the general mediant inequality to these two fractions.

% -------------------------------------------------------
% YOUR PROOF HERE
% -------------------------------------------------------

\end{proof}

\begin{remark*}[Proof shape]
This is a direct corollary of the general mediant inequality.
\end{remark*}

\begin{remark*}[Dependencies]
Uses Theorem~10.49.
\end{remark*}


\subsection*{Lemma 10.70 — Mediant Gap Lemma}
\label{prf:grep-mediant-gap-lemma}

\begin{proof}
\Claim If
\[
\frac{a}{b} < \frac{c}{d},
\]
then the mediant
\[
\frac{a+c}{b+d}
\]
lies between them, and the gaps satisfy
\[
\frac{a+c}{b+d} - \frac{a}{b} = \frac{1}{b(b+d)}
\]
and
\[
\frac{c}{d} - \frac{a+c}{b+d} = \frac{1}{d(b+d)}.
\]

\textbf{Given.} Two ordered rational numbers with positive denominators.

\textbf{Goal.} Compute the two subgap lengths created by inserting the mediant.

\textbf{Strategy.} First establish that the mediant lies between the endpoints.  Then subtract fractions directly and simplify.

% -------------------------------------------------------
% YOUR PROOF HERE
% -------------------------------------------------------

\end{proof}

\begin{remark*}[Proof shape]
This is an explicit computation proof.  The interval-refinement content comes
from exact arithmetic rather than from an abstract limit argument.
\end{remark*}

\begin{remark*}[Dependencies]
Uses Theorem~10.49 and routine fraction subtraction.
\end{remark*}


\subsection*{Theorem 10.54 — Stern--Brocot enumeration}
\label{prf:grep-stern-brocot-enumeration}

\begin{proof}
\Claim Every positive rational number appears exactly once in reduced form in the Stern--Brocot tree.

\textbf{Given.} The Stern--Brocot tree generated from $0/1$ and $1/0$ by repeated mediants.

\textbf{Goal.} Show both existence and uniqueness of appearance for every positive rational.

\textbf{Strategy.} Separate the proof into two parts: (1) every positive rational can be reached by a finite mediant-search path; (2) no reduced rational appears twice.

% -------------------------------------------------------
% YOUR PROOF HERE
% -------------------------------------------------------

\end{proof}

\begin{remark*}[Proof shape]
This is a two-part structural proof.  The existence part is usually handled by a
finite search or Euclidean-algorithm argument; the uniqueness part relies on the
determinant/neighbor structure of Stern--Brocot intervals.
\end{remark*}

\begin{remark*}[Dependencies]
Uses the mediant inequality, Farey/Stern--Brocot determinant phenomena, and
basic properties of reduced fractions.
\end{remark*}


\subsection*{Theorem 10.72 — Rational approximation}
\label{prf:grep-rational-approximation}

\begin{proof}
\Claim For every real number $x$ and every $\varepsilon>0$, there exists $r \in \mathbb{Q}$ such that
\[
|x-r| < \varepsilon.
\]

\textbf{Given.} A real number $x$ and a positive tolerance $\varepsilon$.

\textbf{Goal.} Find a rational number within $\varepsilon$ of $x$.

\textbf{Strategy.} Apply density of $\mathbb{Q}$ to the interval $(x-\varepsilon, x+\varepsilon)$.

% -------------------------------------------------------
% YOUR PROOF HERE
% -------------------------------------------------------

\end{proof}

\begin{remark*}[Proof shape]
This is a one-step application of density.  The metric statement is obtained from
the order statement by choosing a symmetric interval around $x$.
\end{remark*}

\begin{remark*}[Dependencies]
Uses Theorem~10.44.
\end{remark*}


\subsection*{Theorem 10.76 — Dirichlet's Approximation Theorem}
\label{prf:grep-dirichlet-approximation}

\begin{proof}
\Claim For every real number $x$ there exist infinitely many rational numbers $p/q$ with $p \in \mathbb{Z}$ and $q \in \mathbb{N}$ such that
\[
\left|x-\frac{p}{q}\right| < \frac{1}{q^2}.
\]

\textbf{Given.} A real number $x$.

\textbf{Goal.} Produce infinitely many rational approximations satisfying the quadratic error bound.

\textbf{Strategy.} Use the pigeonhole principle on the fractional parts of $x,2x,\dots,(N+1)x$, derive an approximation of size $1/(qN)$, and then turn this into the desired $1/q^2$ bound for infinitely many choices.

% -------------------------------------------------------
% YOUR PROOF HERE
% -------------------------------------------------------

\end{proof}

\begin{remark*}[Proof shape]
This is a combinatorial approximation proof driven by the pigeonhole principle.
The infinite-many conclusion usually comes from applying the finite-$N$ argument
for arbitrarily large $N$.
\end{remark*}

\begin{remark*}[Dependencies]
Uses the pigeonhole principle, properties of fractional parts, and elementary
integer arithmetic.  It does not depend on later Cauchy- or completeness-based
results.
\end{remark*}


\subsection*{Theorem 10.91 — There is no rational number whose square is $2$}
\label{prf:grep-sqrt-two-irrational}

\begin{proof}
\Claim There is no rational number whose square is $2$.

\textbf{Given.} The assumption, for contradiction, that $\sqrt{2}=p/q$ for integers $p,q$ in lowest terms.

\textbf{Goal.} Derive a contradiction.

\textbf{Strategy.} Square the equation, prove that $p$ is even, then prove that $q$ is even, contradicting lowest terms.

% -------------------------------------------------------
% YOUR PROOF HERE
% -------------------------------------------------------

\end{proof}

\begin{remark*}[Proof shape]
This is the classical parity contradiction proof of the irrationality of
$\sqrt{2}$.
\end{remark*}

\begin{remark*}[Dependencies]
Uses reduced-form representatives of rational numbers and elementary parity
facts about integers.
\end{remark*}


\subsection*{Theorem 10.85 — Every convergent sequence in $\mathbb{Q}$ is a Cauchy sequence}
\label{prf:grep-convergent-implies-cauchy}

\begin{proof}
\Claim Every convergent sequence in $\mathbb{Q}$ is a Cauchy sequence.

\textbf{Given.} A sequence $(x_n)$ in $\mathbb{Q}$ converging to some $x \in \mathbb{Q}$.

\textbf{Goal.} Show that for every $\varepsilon>0$ there exists $N$ such that $m,n\ge N$ implies $|x_m-x_n|<\varepsilon$.

\textbf{Strategy.} Use the definition of convergence with tolerance $\varepsilon/2$ and apply the triangle inequality.

% -------------------------------------------------------
% YOUR PROOF HERE
% -------------------------------------------------------

\end{proof}

\begin{remark*}[Proof shape]
This is the standard $\varepsilon/2$ triangle-inequality argument.
\end{remark*}

\begin{remark*}[Dependencies]
Uses the definition of convergence in $\mathbb{Q}$, the definition of Cauchy
sequence, and the triangle inequality.
\end{remark*}


\subsection*{Theorem 10.88 — There exist Cauchy sequences in $\mathbb{Q}$ that do not converge to a rational number}
\label{prf:grep-cauchy-not-convergent-in-q}

\begin{proof}
\Claim There exist Cauchy sequences in $\mathbb{Q}$ that do not converge to a rational number.

\textbf{Given.} A candidate sequence of rational approximations to $\sqrt{2}$, for example decimal truncations or another explicit approximating sequence.

\textbf{Goal.} Show that the sequence is Cauchy but has no rational limit.

\textbf{Strategy.} Split the proof into two parts: (1) verify the sequence is Cauchy; (2) prove that if it converged in $\mathbb{Q}$, its limit would have square $2$, contradicting Theorem~10.91.

% -------------------------------------------------------
% YOUR PROOF HERE
% -------------------------------------------------------

\end{proof}

\begin{remark*}[Proof shape]
This is the first genuine incompleteness proof in the chapter.  The Cauchy part
is analytic; the nonconvergence part is arithmetic and should terminate in the
irrationality of $\sqrt{2}$.
\end{remark*}

\begin{remark*}[Dependencies]
Uses Theorem~10.91 and the definition of Cauchy sequence.  Depending on the
chosen example, may also use a monotone-bound argument or direct decimal-error
estimates.
\end{remark*}


\subsection*{Theorem 10.96 — Stern--Brocot interval refinement}
\label{prf:grep-stern-brocot-interval-refinement}

\begin{proof}
\Claim Let $x>0$ be a real number.  Suppose one begins with two rational numbers
\[
\frac{a}{b} < x < \frac{c}{d}, \qquad b,d>0,
\]
and repeatedly inserts the mediant
\[
\frac{a+c}{b+d}.
\]
At each step, replace one endpoint by the mediant so that the new interval still contains $x$.  Then this process generates a nested sequence of rational intervals containing $x$.

\textbf{Given.} A target real number $x>0$ and an initial rational interval containing it.

\textbf{Goal.} Show that the iterative mediant-selection process yields nested rational intervals, each still containing $x$.

\textbf{Strategy.} Apply the mediant inequality at each step.  Since the mediant lies between the endpoints, replacing one endpoint preserves containment and makes the interval smaller.

% -------------------------------------------------------
% YOUR PROOF HERE
% -------------------------------------------------------

\end{proof}

\begin{remark*}[Proof shape]
This is an induction-on-steps proof for an iterative construction.
Each stage is justified by the mediant inequality.
\end{remark*}

\begin{remark*}[Dependencies]
Uses Theorem~10.49 and the order properties of real numbers.
\end{remark*}


\subsection*{Theorem 10.103 — Stern--Brocot Limit Theorem}
\label{prf:grep-stern-brocot-limit-theorem}

\begin{proof}
\Claim Every irrational number $x>0$ determines a unique infinite path in the Stern--Brocot tree.  The sequence of rational numbers encountered along this path converges to $x$.
Conversely, every infinite path in the Stern--Brocot tree determines a unique positive irrational number as its limit.

\textbf{Given.} An irrational number $x>0$, or alternatively an infinite path in the Stern--Brocot tree.

\textbf{Goal.} Prove both directions: existence/uniqueness of the path from $x$, and existence/uniqueness of the limit from the path.

\textbf{Strategy.} For the forward direction, repeatedly refine Stern--Brocot intervals according to whether $x$ lies left or right of the mediant.  For the reverse direction, show that an infinite path yields nested rational intervals whose lengths tend to $0$, and use this to define a unique real limit point.

% -------------------------------------------------------
% YOUR PROOF HERE
% -------------------------------------------------------

\end{proof}

\begin{remark*}[Proof shape]
This is the capstone structural proof of the chapter.  It is naturally divided
into two directions and is best approached through nested intervals whose
lengths shrink to zero.
\end{remark*}

\begin{remark*}[Dependencies]
Uses Theorem~10.96, the mediant-gap structure (especially Lemma~10.70), and a
nested-interval or completeness argument in $\mathbb{R}$.
\end{remark*}


% =========================================================
% Proof stubs — Stern–Brocot Limit Theory
% Four files; each would live at:
%   volume-ii/rationals/proofs/notes/<filename>.tex
%
% File A: sb-denom-growth.tex
% File B: sb-length-zero.tex
% File C: sb-unique-point.tex
% File D: sb-path-limit.tex
% File E: sb-limit-irrational.tex
% =========================================================


% =========================================================
% FILE A: sb-denom-growth.tex
% Proof: Lemma — Denominator growth on infinite SB paths
% Source: volume-ii/rationals/notes/<topic>.tex
% =========================================================

\subsection*{Lemma --- Denominator growth on infinite Stern--Brocot paths}
\label{prf:sb-denom-growth}

\begin{remark*}[Return]
$\leftarrow$ Back to Lemma~\ref{lem:sb-denom-growth} (Denominator growth) in Notes.
\end{remark*}

\begin{proof}
\Claim Along any infinite path in the Stern--Brocot tree, the denominators
of the successive endpoint fractions form a strictly increasing sequence
tending to $\infty$.

\Given A path in the Stern--Brocot tree.  At each level $n$ the path
occupies an interval $[a_n/b_n,\; c_n/d_n]$.  By the Farey determinant
condition (Farey Neighbor Theorem, Theorem~10.61), the two endpoints satisfy
$b_n c_n - a_n d_n = 1$ at every level.

\Goal To show $b_{n+1} > b_n$ (and symmetrically $d_{n+1} > d_n$) for
every $n$, and hence $b_n \to \infty$.

\Strategy We proceed by induction on the level $n$ of the path.

\textbf{Base case} ($n = 0$). The tree is initialised from $0/1$ and $1/0$
(the sentinel fractions). Their denominators are $1$ and $0$. When the
first mediant $1/1$ is inserted, the new interval is either $[0/1, 1/1]$
with denominators $1, 1$, or $[1/1, 1/0]$ with denominators $1, 0$.
In both cases the denominator of the chosen boundary fraction is $b_0 = 1$,
which is positive. \checkbox

\textbf{Inductive step.} Assume the interval at level $n$ is
$[a_n/b_n,\; c_n/d_n]$ with $b_n, d_n \ge 1$.

\ByDef the mediant-insertion rule, the new fraction inserted at level
$n+1$ is
\[
  \frac{a_n + c_n}{b_n + d_n}.
\]
One endpoint of the new interval is this mediant, which has denominator
$b_n + d_n$.

\ByAss $b_n \ge 1$ and $d_n \ge 1$, we have
\[
  b_n + d_n > b_n \quad \text{and} \quad b_n + d_n > d_n.
\]
\Therefore the denominator of the new endpoint strictly exceeds both
$b_n$ and $d_n$. The other endpoint retains the denominator it had
at level $n$.  In either case the minimum denominator on the boundary
strictly increases. \checkbox

\Hence by induction every denominator encountered along the path is
strictly larger than its predecessor.  Since the denominators are positive
integers, a strictly increasing sequence of positive integers is unbounded.
\Therefore $b_n \to \infty$. \AsReq
\end{proof}

\begin{remark*}[Proof shape]
Induction on depth, with the key computation being $b_n + d_n > b_n$.
The unboundedness conclusion is immediate once strict monotonicity of a
positive-integer sequence is in hand.
\end{remark*}

\begin{remark*}[Dependencies]
Farey Neighbor Theorem (Theorem~10.61) --- guarantees the determinant
condition at each level, confirming the endpoints are genuine Farey
neighbors. Positivity of natural numbers (Corollary~10.26) --- ensures
denominators are at least $1$ after the first step, making the strict
inequality $b_n + d_n > b_n$ valid.
\end{remark*}


% =========================================================
% FILE B: sb-length-zero.tex
% Proof: Lemma — SB interval lengths shrink to zero
% Source: volume-ii/rationals/notes/<topic>.tex
% =========================================================

\subsection*{Lemma --- Stern--Brocot interval lengths shrink to zero}
\label{prf:sb-length-zero}

\begin{remark*}[Return]
$\leftarrow$ Back to Lemma~\ref{lem:sb-length-zero} (SB lengths to zero) in Notes.
\end{remark*}

\begin{proof}
\Claim Let $(I_n)$ be the nested intervals from the Stern--Brocot
refinement process, with $I_n = [a_n/b_n,\; c_n/d_n]$ and
$\ell_n = c_n/d_n - a_n/b_n$. Then $\ell_n \to 0$.

\Given At each level $n$, the two endpoints $a_n/b_n$ and $c_n/d_n$
are Farey neighbors (the determinant condition is preserved by mediant
insertion, as verified in the proof of the Farey Neighbor Theorem).

\Goal To show $\forall \varepsilon > 0\ \exists N\ \forall n \ge N,\ \ell_n < \varepsilon$.

\Strategy Apply the Gap Between Farey Neighbors identity, then use the
denominator growth lemma.

\ByThm{Gap Between Farey Neighbors} (Theorem~10.65), since $a_n/b_n$
and $c_n/d_n$ are Farey neighbors at every level,
\[
  \ell_n = \frac{c_n}{d_n} - \frac{a_n}{b_n} = \frac{1}{b_n d_n}.
\]

\ByThm{Denominator growth on infinite Stern--Brocot paths}
(Lemma~\ref{lem:sb-denom-growth}), $b_n \to \infty$ and $d_n \to \infty$.

\Let $\varepsilon > 0$ be arbitrary.
\ByThm{Archimedean property of $\mathbb{Q}$} (Theorem~10.42), there
exists $N \in \mathbb{N}$ such that $b_N d_N > 1/\varepsilon$.

For all $n \ge N$, the strict monotonicity of denominators gives
$b_n \ge b_N$ and $d_n \ge d_N$, so
\[
  \ell_n = \frac{1}{b_n d_n} \le \frac{1}{b_N d_N} < \varepsilon.
\]

\Hence $\ell_n \to 0$. \AsReq
\end{proof}

\begin{remark*}[Proof shape]
A direct application of two previously established results: the exact
gap formula and the denominator growth lemma.  The Archimedean property
converts the asymptotic statement into a concrete $N$.
\end{remark*}

\begin{remark*}[Dependencies]
Gap Between Farey Neighbors (Theorem~10.65) --- provides the identity
$\ell_n = 1/(b_n d_n)$. Denominator growth on infinite Stern--Brocot
paths (Lemma~\ref{lem:sb-denom-growth}) --- supplies $b_n, d_n \to \infty$.
Archimedean property of $\mathbb{Q}$ (Theorem~10.42) --- used to find
the index $N$.
\end{remark*}


% =========================================================
% FILE C: sb-unique-point.tex
% Proof: Proposition — Nested SB intervals contain unique point
% Source: volume-ii/rationals/notes/<topic>.tex
% =========================================================

\subsection*{Proposition --- Nested Stern--Brocot intervals contain a unique limit point}
\label{prf:sb-unique-point}

\begin{remark*}[Return]
$\leftarrow$ Back to Proposition~\ref{prop:sb-unique-point} (Unique limit point) in Notes.
\end{remark*}

\begin{proof}
\Claim Let $(I_n)$ with $I_n = [a_n/b_n,\; c_n/d_n]$ be the nested
closed rational intervals from the Stern--Brocot refinement process.
Then $\bigcap_{n=0}^{\infty} I_n$ contains exactly one point.

\Given The intervals are nested: $I_0 \supseteq I_1 \supseteq I_2 \supseteq \cdots$
by construction (at each step one endpoint is replaced by the mediant,
which lies strictly inside the current interval). The lengths satisfy
$\ell_n \to 0$ by Lemma~\ref{lem:sb-length-zero}.

\Goal Existence: $\bigcap I_n \ne \varnothing$.
Uniqueness: $|\bigcap I_n| = 1$.

\Strategy Apply the Monotone Convergence Theorem in $\mathbb{R}$ to both
endpoint sequences; then use $\ell_n \to 0$ to pin them to the same limit.

\textbf{Left endpoints.} The sequence $(a_n/b_n)$ is increasing
(each left endpoint is either unchanged or replaced by a larger mediant)
and bounded above by $c_0/d_0$. \ByThm{Monotone Convergence Theorem in
$\mathbb{R}$}, the supremum $\alpha = \sup_n a_n/b_n$ exists in
$\mathbb{R}$ and $a_n/b_n \to \alpha$.

\textbf{Right endpoints.} The sequence $(c_n/d_n)$ is decreasing
and bounded below by $a_0/b_0$.  By the same theorem, the infimum
$\beta = \inf_n c_n/d_n$ exists in $\mathbb{R}$ and $c_n/d_n \to \beta$.

\Therefore $\alpha \le \beta$ (since $a_n/b_n \le c_n/d_n$ for all $n$).
But
\[
  0 \le \beta - \alpha \le \ell_n = \frac{1}{b_n d_n} \to 0,
\]
so $\alpha = \beta$.

\Let $\xi = \alpha = \beta$. Then $a_n/b_n \le \xi \le c_n/d_n$ for
all $n$, so $\xi \in I_n$ for all $n$.
\Therefore $\xi \in \bigcap I_n$, establishing existence.

\textbf{Uniqueness.} Suppose $\xi, \xi' \in \bigcap I_n$.
Then for every $n$,
\[
  |\xi - \xi'| \le \ell_n = \frac{1}{b_n d_n}.
\]
Since $\ell_n \to 0$, we have $|\xi - \xi'| = 0$, so $\xi = \xi'$.

\Hence $\bigcap I_n = \{\xi\}$. \AsReq
\end{proof}

\begin{remark*}[Proof shape]
A nested-intervals argument in $\mathbb{R}$: the left endpoints form a
bounded increasing sequence and the right endpoints form a bounded decreasing
sequence; both converge; the limit length forces the two limits to coincide.
Uniqueness is a direct consequence of the length tending to zero.
\end{remark*}

\begin{remark*}[Dependencies]
Monotone Convergence Theorem in $\mathbb{R}$ (from the bounding/sequences
chapter) --- the load-bearing tool; requires completeness of $\mathbb{R}$,
which is why the proof cannot be carried out inside $\mathbb{Q}$.
Stern--Brocot interval lengths shrink to zero (Lemma~\ref{lem:sb-length-zero})
--- supplies $\ell_n \to 0$.
\end{remark*}


% =========================================================
% FILE D: sb-path-limit.tex
% Proof: Corollary — Every infinite SB path converges
% Source: volume-ii/rationals/notes/<topic>.tex
% =========================================================

\subsection*{Corollary --- Every infinite Stern--Brocot path converges}
\label{prf:sb-path-limit}

\begin{remark*}[Return]
$\leftarrow$ Back to Corollary~\ref{cor:sb-path-limit} (SB path converges) in Notes.
\end{remark*}

\begin{proof}
\Claim Both endpoint sequences $(a_n/b_n)$ and $(c_n/d_n)$ of the
Stern--Brocot refinement process targeting $x > 0$ converge to $x$.

\Given The process constructs $I_n = [a_n/b_n, c_n/d_n]$ with $x \in I_n$
for all $n$ (by the selection rule: at each step the endpoint that does not
separate $x$ from the mediant is retained, so $x$ remains in the interval).

\Strategy Apply Proposition~\ref{prop:sb-unique-point}.

\ByThm{Nested Stern--Brocot intervals contain a unique limit point}
(Proposition~\ref{prop:sb-unique-point}), there is a unique $\xi$ with
$\xi \in \bigcap I_n$ and $a_n/b_n, c_n/d_n \to \xi$.

\ByAss $x \in I_n$ for all $n$, so $x \in \bigcap I_n$.
By the uniqueness from Proposition~\ref{prop:sb-unique-point}, $\xi = x$.

\Hence $a_n/b_n \to x$ and $c_n/d_n \to x$. \AsReq
\end{proof}

\begin{remark*}[Proof shape]
A one-step corollary: the unique point in $\bigcap I_n$ is $\xi$;
$x$ also lies in every $I_n$; therefore $\xi = x$.
\end{remark*}

\begin{remark*}[Dependencies]
Nested Stern--Brocot intervals contain a unique limit point
(Proposition~\ref{prop:sb-unique-point}).
Construction rule for the Stern--Brocot refinement process (the selection
step that keeps $x$ in the interval at every stage).
\end{remark*}


% =========================================================
% FILE E: sb-limit-irrational.tex
% Proof: Corollary — Infinite paths converge to irrationals
% Source: volume-ii/rationals/notes/<topic>.tex
% =========================================================

\subsection*{Corollary --- Infinite Stern--Brocot paths converge to irrationals}
\label{prf:sb-limit-irrational}

\begin{remark*}[Return]
$\leftarrow$ Back to Corollary~\ref{cor:sb-limit-irrational} (Irrational limits) in Notes.
\end{remark*}

\begin{proof}
\Claim A positive real number $x$ is rational if and only if it
corresponds to a finite path in the Stern--Brocot tree.

\Given The Stern--Brocot tree is generated from $0/1$ and $1/0$ by
repeated mediant insertion. At each step the process asks: is $x$ equal
to, less than, or greater than the current mediant?

\Strategy We prove both implications separately.

\medskip
\noindent$(\Rightarrow)$\ \textit{Rational implies finite path.}

\Given $x = p/q \in \mathbb{Q}_{>0}$ in reduced form.

The search process compares $x$ to successive mediants.  At each step
the two boundary fractions $a_n/b_n$ and $c_n/d_n$ satisfy $b_n c_n -
a_n d_n = 1$ (Farey Neighbor Theorem, Theorem~10.61) and the mediant
has denominator $b_n + d_n$.

Suppose $x$ is never equal to a mediant. Then at every level the
inequality is strict and both $b_n$ and $d_n$ are at most $q - 1$
(otherwise a fraction with denominator $\le q$ equal to $x$ would have
been produced, contradicting the assumption that $x = p/q$ is in reduced
form with denominator exactly $q$). But Lemma~\ref{lem:sb-denom-growth}
says $b_n \to \infty$, contradiction.

\Therefore the process must encounter $x$ as a mediant at some finite
level $N$, and the path terminates there. \checkbox

\medskip
\noindent$(\Leftarrow)$\ \textit{Finite path implies rational.}

\Given The path terminates at level $N$, meaning $x$ equals the mediant
$(a_N + c_N)/(b_N + d_N)$.

\ByDef a mediant is a rational number (sum of integers over sum of
integers).
\Therefore $x \in \mathbb{Q}$. \checkbox

\medskip
Taking the contrapositive of the $(\Rightarrow)$ direction:
if $x$ is irrational then the path is infinite.
Together with $(\Leftarrow)$ this gives the equivalence. \AsReq
\end{proof}

\begin{remark*}[Proof shape]
A biconditional proof. The $(\Leftarrow)$ direction is trivial (a mediant
is visibly rational). The $(\Rightarrow)$ direction is the substantive half:
it proceeds by contradiction, using denominator growth to force $x$ to
appear as a mediant within a bounded number of steps.
\end{remark*}

\begin{remark*}[Dependencies]
Farey Neighbor Theorem (Theorem~10.61) --- the determinant condition
$b_n c_n - a_n d_n = 1$ controls what denominators can appear. Denominator
growth on infinite Stern--Brocot paths (Lemma~\ref{lem:sb-denom-growth})
--- the engine of the $(\Rightarrow)$ contradiction. Definition of reduced
form for rational numbers (Definition~10.8 and Remark~10.15).
\end{remark*}